\documentclass[10pt, letterpaper]{article}

% Packages:
\usepackage[
    ignoreheadfoot, % set margins without considering header and footer
    top=0.5 in, % seperation between body and page edge from the top
    bottom=0.5 in, % seperation between body and page edge from the bottom
    left=0.8 in, % seperation between body and page edge from the left
    right=0.8 in, % seperation between body and page edge from the right
    footskip=0.5 in, % seperation between body and footer
    % showframe % for debugging 
]{geometry} % for adjusting page geometry

% \usepackage{hyperref}

\usepackage{titlesec} % for customizing section titles
\usepackage{tabularx} % for making tables with fixed width columns
\usepackage{array} % tabularx requires this
\usepackage[dvipsnames]{xcolor} % for coloring text
\definecolor{primaryColor}{RGB}{0, 0, 0} % define primary color
\usepackage{enumitem} % for customizing lists
\usepackage{fontawesome5} % for using icons
\usepackage{amsmath} % for math
\usepackage{multicol}
\usepackage[
    pdftitle={Jose Tupayachi - Resume},
    pdfauthor={Jose Tupayachi},
    pdfcreator={LaTeX with RenderCV},
    colorlinks=true,
    urlcolor=primaryColor
]{hyperref} % for links, metadata and bookmarks
\usepackage[pscoord]{eso-pic} % for floating text on the page
\usepackage{calc} % for calculating lengths
\usepackage{bookmark} % for bookmarks
\usepackage{lastpage} % for getting the total number of pages
\usepackage{changepage} % for one column entries (adjustwidth environment)
\usepackage{paracol} % for two and three column entries
\usepackage{ifthen} % for conditional statements
\usepackage{needspace} % for avoiding page brake right after the section title
\usepackage{iftex} % check if engine is pdflatex, xetex or luatex

% Ensure that generate pdf is machine readable/ATS parsable:
\ifPDFTeX
    \input{glyphtounicode}
    \pdfgentounicode=1
    \usepackage[T1]{fontenc}
    \usepackage[utf8]{inputenc}
    \usepackage{lmodern}
\fi

\usepackage{charter}

% Some settings:
\raggedright
\AtBeginEnvironment{adjustwidth}{\partopsep0pt} % remove space before adjustwidth environment
\pagestyle{empty} % no header or footer
\setcounter{secnumdepth}{0} % no section numbering
\setlength{\parindent}{0pt} % no indentation
\setlength{\topskip}{0pt} % no top skip
\setlength{\columnsep}{0.15cm} % set column seperation
\pagenumbering{gobble} % no page numbering

\titleformat{\section}{\needspace{4\baselineskip}\bfseries\large}{}{0pt}{}[\vspace{1pt}\titlerule]

\titlespacing{\section}{
    % left space:
    -1pt
}{
    % top space:
    0.3 cm
}{
    % bottom space:
    0.2 cm
} % section title spacing

\renewcommand\labelitemi{$\vcenter{\hbox{\small$\bullet$}}$} % custom bullet points
\newenvironment{highlights}{
    \begin{itemize}[
        topsep=0.10 cm,
        parsep=0.10 cm,
        partopsep=0pt,
        itemsep=0pt,
        leftmargin=0 cm + 10pt
    ]
}{
    \end{itemize}
} % new environment for highlights


\newenvironment{highlightsforbulletentries}{
    \begin{itemize}[
        topsep=0.10 cm,
        parsep=0.10 cm,
        partopsep=0pt,
        itemsep=0pt,
        leftmargin=10pt
    ]
}{
    \end{itemize}
} % new environment for highlights for bullet entries

\newenvironment{onecolentry}{
    \begin{adjustwidth}{
        0 cm + 0.00001 cm
    }{
        0 cm + 0.00001 cm
    }
}{
    \end{adjustwidth}
} % new environment for one column entries

\newenvironment{twocolentry}[2][]{
    \onecolentry
    \def\secondColumn{#2}
    \setcolumnwidth{\fill, 4.5 cm}
    \begin{paracol}{2}
}{
    \switchcolumn \raggedleft \secondColumn
    \end{paracol}
    \endonecolentry
} % new environment for two column entries

\newenvironment{threecolentry}[3][]{
    \onecolentry
    \def\thirdColumn{#3}
    \setcolumnwidth{, \fill, 4.5 cm}
    \begin{paracol}{3}
    {\raggedright #2} \switchcolumn
}{
    \switchcolumn \raggedleft \thirdColumn
    \end{paracol}
    \endonecolentry
} % new environment for three column entries

\newenvironment{header}{
    \setlength{\topsep}{0pt}\par\kern\topsep\centering\linespread{1.5}
}{
    \par\kern\topsep
} % new environment for the header

\newcommand{\placelastupdatedtext}{% \placetextbox{<horizontal pos>}{<vertical pos>}{<stuff>}
  \AddToShipoutPictureFG*{% Add <stuff> to current page foreground
    \put(
        \LenToUnit{\paperwidth-2 cm-0 cm+0.05cm},
        \LenToUnit{\paperheight-1.0 cm}
    ){\vtop{{\null}\makebox[0pt][c]{
        \small\color{gray}\textit{Last updated in July 2026}\hspace{\widthof{Last updated in July 2026}}
    }}}%
  }%
}%

% save the original href command in a new command:
\let\hrefWithoutArrow\href

% new command for external links:


\begin{document}
    \newcommand{\AND}{\unskip
        \cleaders\copy\ANDbox\hskip\wd\ANDbox
        \ignorespaces
    }
    \newsavebox\ANDbox
    \sbox\ANDbox{$|$}

    \begin{header}
        \fontsize{22 pt}{22 pt}\selectfont \textbf{Jose Tupayachi}

        \vspace{3 pt}

        \normalsize
        Data and AI Engineer

        \vspace{3 pt}

        \small
        \mbox{Knoxville, TN}%
        \kern 5.0 pt%
        \AND%
        \kern 5.0 pt%
        \mbox{\hrefWithoutArrow{mailto:jtupayac@vols.utk.edu}{jtupayac@vols.utk.edu}}%
        \kern 5.0 pt%
        \AND%
        \kern 5.0 pt%
        \mbox{\hrefWithoutArrow{mailto:tupayachisja@ornl.gov}{tupayachisja@ornl.gov}}%
        \kern 5.0 pt%
        \AND%
        \kern 5.0 pt%
        \mbox{\hrefWithoutArrow{tel:+16613655289}{661-365-5289}}%
        \kern 5.0 pt%
        \AND%
        \kern 5.0 pt%
        \mbox{\hrefWithoutArrow{https://jtupayachi.github.io/}{jtupayachi.github.io}}%
    \end{header}
% ----------------------------
\section{SUMMARY}
% ----------------------------
\small

% Jose Tupayachi is a data and AI engineer with full stack expertise: Scalable ETL pipelines, big data platforms (Spark, Databricks), cloud data warehouses, and advanced machine learning.
% He designs and deploys end to end data pipelines (Ingestion and transformation to predictive modeling and decision support) integrating large language models, retrieval-augmented generation (RAG), and multimodal AI to turn complex datasets into actionable insights.
% His applied work also spans into AI for multimodal information: This includes materials sciences (SEM analysis), AI powered freight transportation digital twins, geospatial analytics, and Android-iOS platforms.
% He has published in \textit{Computers and Electrical Engineering}, \textit{Health Informatics Journal}, and has filed IP disclosures at Oak Ridge National Laboratory.

% Jose Tupayachi is a Data and AI Engineer specializing in scalable data engineering, machine learning, and multimodal AI. His expertise includes building end-to-end data platforms using Apache Spark, Databricks, cloud data warehouses, and modern MLOps practices, while integrating large language models (LLMs), retrieval-augmented generation (RAG), and multimodal AI into production workflows. His research and applied work span materials science (SEM image analysis), AI-powered freight transportation digital twins, geospatial analytics, and mobile application development. He has published in \textit{Computers \& Electrical Engineering} and the \textit{Health Informatics Journal}, and has filed intellectual property disclosures at Oak Ridge National Laboratory.


I specialize in scalable ETL, distributed computing, machine learning, and multimodal AI. Throughout my career, I have deployed pipelines on a various cloud platforms designing data models and datalake architectures. I engineer end-to-end, agentic data workflows, from ingestion, transformation, and orchestration to predictive modeling and decision support, integrating large language models, retrieval-augmented generation, vision transformers, and agentic AI to turn complex unstructured information into actionable insights. My applied work spans finance, materials science, freight transportation digital twins, geospatial analytics, and cross-platform mobile development. I have published in \textit{Computers \& Electrical Engineering} and the \textit{Health Informatics Journal}, and filed intellectual property disclosures at Oak Ridge National Laboratory.


% ----------------------------
\section{EDUCATION}
% ----------------------------
\small

\begin{itemize}[topsep=0pt, parsep=0pt, partopsep=0pt, itemsep=3pt, leftmargin=10pt]
    \item \textbf{Ph.D.}, Industrial Engineering, University of Tennessee, Knoxville, May 2024 - Present \\
          \textit{Advisor: Dr.\ Xueping Li} | \textit{Dr.\ Xiao-Ying Yu}
    \item \textbf{M.S.}, Industrial Engineering, University of Tennessee, Knoxville, Aug 2022 - May 2024
    \item \textbf{B.S.}, Industrial Engineering, Pontifical Catholic University of Peru, March 2014 - Aug 2020
\end{itemize}

% ----------------------------
\section{EXPERIENCE}
% ----------------------------
\small

\noindent\textbf{Developer}, Oak Ridge National Laboratory / ORISE \hfill Mar 2025 - Present\\
\textit{Data Engineer developing multimodal pipelines and AI systems for materials science, fusion energy, and geospatial data under DOE-funded programs.}
\begin{highlights}
    \item Built Vision Transformer (ViT) and fine-tuned vision models within a multi-pathway hybrid architecture to deploy an \textit{Automated Quality Assurance of Meteorological Station Data for Atmospheric Dispersion Modeling (CAPARS)} to classify wind speed and direction across distributed meteorological networks.
    \item Developed agentic, traceable multimodal AI pipelines with LangChain and LangGraph, integrating large language models (LLMs) and retrieval-augmented generation (RAG) to automate expert literature review and new research hyphotesis workflows.
    \item Fine-tuned and deployed local open models (Qwen, Llama) served via Ollama for autonomous literature mining and knowledge extraction.
    \item Created explainable ML models for extreme weather and environmental collection towers using large-scale atmospheric and geospatial datasets.
    \item Developed geospatial AI dashboard to enable natural-language querying, knowledge-graph reasoning, and interactive analytics over spatial datasets, filed as ORNL intellectual property disclosure.
    % \item Built LLM data analysis decision support and automated dashboards, .
\end{highlights}


\vspace{0.2cm}
\noindent\textbf{Graduate Research Assistant}, University of Tennessee, Knoxville \hfill Jul 2024 - Present\\
\textit{Researcher in AI-driven transportation systems, digital twins, and multimodal machine learning for mobility, healthcare, and infrastructure applications.}
\begin{highlights}
    \item Lead development of AI-powered freight transportation digital twins under the \$3M DOE ARPA-E RECOIL project, delivering real-time spatiotemporal prediction and system optimization.
    \item Designed and implemented deep learning architectures (MLP, CNN, LSTM, transformer-based models) for EV charging-demand forecasting and spatiotemporal prediction.
    \item Built agentic LLM workflows and RAG chatbots for domain-specific trade-off (what-if) analysis and deployed to production for non-technical decision-makers.
    \item Delivered multimodal ML systems for cross-platform mobile health apps (SmartShots, ACTAPP), published to the App Store and Google Play.
    \item Mentored junior students in model development, GitHub and GNU parallel to speed up execution times.
\end{highlights}

\vspace{0.2cm}
\noindent\textbf{Teaching Assistant - IE406/408: Simulation}, University of Tennessee, Knoxville \hfill Jan 2022 - May 2022
\begin{highlights}
    \item Supported instruction in discrete-event simulation and led lab sessions.
    \item Assisted with AnyLogic modeling assignments and graded coursework.
\end{highlights}

\vspace{0.2cm}
\noindent\textbf{Teaching Assistant - IE483: Reliability Engineering}, University of Tennessee, Knoxville \hfill Aug 2022 - Dec 2022
\begin{highlights}
    \item Assisted with reliability analysis topics including statistical distributions, failure mode analysis, and remaining useful life prediction methods.
    \item Held office hours.
\end{highlights}

\vspace{0.2cm}
\noindent\textbf{Data Engineer}, Indra \hfill Mar 2022 - Aug 2022\\
\textit{Data engineer supporting enterprise-scale ETL pipelines, distributed data processing, and cloud migration for large-scale analytics systems.}
\begin{highlights}
    \item Built and maintained scalable ETL/ELT pipelines using Python, Scala, and Shell scripting on UNIX/Linux, processing high-volume datasets data in Agile delivery cycles.
    \item Developed distributed data processing systems using Scala/PySpark on Apache Spark, Hadoop, Hive/HQL, and AWS (S3, EMR), migrating legacy Oracle warehouse to a partitioned cloud lakehouse and tuning Spark memory and resource allocation for faster job execution.
    \item Designed dimensional data models for data warehousing and implemented entity-resolution and deduplication logic to unify records across different tables.
    \item Implemented CI/CD deployment automation with Jenkins for automated builds, unit testing, and application deployment (live deployment, rollback and backups).
    \item Implemented a data quality (ad-hoc quality) engine, profiling, and drifting detection through validation frameworks and monitoring/alerting.
\end{highlights}



\vspace{0.2cm}
\noindent\textbf{Business Intelligence Analyst}, Globokas \hfill Jan 2022 - Mar 2022\\
\textit{BI analyst designing data-driven reporting systems, dashboards, and analytical models for business decision support.}
\begin{highlights}
    \item Developed executive dashboards and KPI reporting systems supporting strategic and operational decision-making for multiple business clients.
    \item Integrated heterogeneous data sources (REST APIs, internal systems, external platforms) into unified analytics pipelines.
    \item Translated business requirements into analytical models and actionable insights for stakeholders.
\end{highlights}

\vspace{0.2cm}
\noindent\textbf{Data Analyst}, Enel Distribución \hfill Nov 2020 - Dec 2021\\
\textit{Data analyst building predictive models, customer segmentation systems, and automated reporting tools for energy operations.}
\begin{highlights}
    \item Applied unsupervised learning (clustering) for large-scale customer segmentation to improve payment-collection targeting and recovery.
    \item Built interactive Power BI and Tableau dashboards for operational monitoring and KPI reporting.
    \item Managed SQL/T-SQL databases and Salesforce integrations API queries.
    \item Developed a PyQt desktop application to automate invoice verification and special cases processing.
\end{highlights}

\vspace{0.2cm}
\noindent\textbf{Developer - Portfolio Valuation}, GPS Management \hfill Aug 2020 - Nov 2020
\begin{highlights}
    \item Implemented machine learning models for non-performing loan (NPL) recovery and portfolio valuation.
    \item Automated scoring workflows, reducing manual processing time.
\end{highlights}

\vspace{0.2cm}
\noindent\textbf{Professional IT Intern}, Superintendencia de Banca, Seguros y AFP \hfill Sep 2020 - Oct 2020
\begin{highlights}
    \item Supported IT infrastructure and regulatory data management.
    \item Assisted in database maintenance and internal reporting for financial supervision operations.
\end{highlights}

\vspace{0.2cm}
\noindent\textbf{IT Intern}, Administradora Clínica Ricardo Palma \hfill Feb 2019 - Aug 2019
\begin{highlights}
    \item Developed and maintained internal healthcare information systems.
    \item Supported electronic health record integration and database administration for clinical operations.
\end{highlights}

% ----------------------------
\section{PROJECTS}
% ----------------------------
\small



\noindent\textbf{Oak Ridge National Laboratory} \hfill Mar 2025 - Present

\noindent\textbf{Role:} Research Fellow. Deployed advanced computer vision, multimodal AI, and scientific computing systems for materials science and engineering applications. Developed Vision Transformer (ViT) architectures, foundation vision models, and AI-assisted characterization pipelines for scanning electron microscopy (SEM) imagery, enabling automated materials analysis, microstructure characterization, defect detection, feature extraction, and accelerated materials discovery workflows. By integrating large language models, multimodal foundation models, retrieval-augmented generation (RAG), and domain-specific scientific knowledge, built intelligent research platforms that transform complex experimental and simulation data into actionable insights for scientists and engineers.

\noindent\textbf{Funding:} DOE Office of Science (SC), Fusion Energy Sciences (FES) - Fusion Materials Sciences Program, DOE Office of Environmental Management (EM), Nuclear Safety Research and Development (NSRD).

\begin{highlights}
    \item Developed AI-agent workflows for materials discovery, scientific knowledge extraction, and literature exploration.
    % \item Built computer vision, Vision Transformer (ViT), and foundation vision models for scanning electron microscopy (SEM) image characterization and materials science applications.
    \item Applied multimodal AI systems to automate microstructure analysis, defect identification, phase characterization, and feature extraction from scanning electron microscopy (SEM) image.
    \item Developed AI-driven characterization pipelines supporting fusion materials research.
    \item Designed explainable machine learning models for extreme weather event prediction using atmospheric, environmental, and meteorological towers.
    \item Developed geospatial AI interfaces enabling natural-language interaction with spatial datasets and  analysis workflows.
    \item Created automated dashboard generation and scientific decision-support systems using large language models, retrieval-augmented generation (RAG) using web technologies.
\end{highlights}


\vspace{0.3cm}
\noindent\textbf{Intellectual Property Disclosures} \textbar\ Oak Ridge National Laboratory \hfill 2025 - Present

\noindent\textit{Automated Geospatial Analysis Interface} - engineered an end-to-end interface driven by large language models, enabling non-expert users to query, process, and visualize spatial datasets through natural language. \textit{IP disclosure filed, ORNL, 2026.}


\vspace{0.3cm}
\noindent\textbf{RECOIL} - Cognitive Freight Transportation Digital Twin \hfill Jul 2024 - March 2026

\noindent\textbf{Role:} Lead AI/Data Engineer. Designed and built an end-to-end cognitive digital twin for large-scale intermodal freight transportation, integrating optimization, deep learning, large language models, and real-time streaming to minimize costs and $CO_2$ emissions.

\noindent\textbf{Funding:} U.S. Department of Energy ARPA-E, Project \#DE-AR0001780.

\begin{highlights}
    \item Designed ontology-guided optimization models integrating GIS-based transport modes (road, rail, waterways) for cost and emissions minimization.
    \item Leveraged Convolutional and Graph Neural Networks to analyze weather, traffic, and demand impacts on EV charging behavior for real-time optimization.
    \item Applied large language models to power a domain-specific chatbot for trade-off analysis between cost, time, and $CO_2$ emissions for non-technical users.
    \item Designed a real-time feedback loop using simulated sensor data and streaming ingestion to continuously optimize intermodal system performance.
    \item Deployed BERT and sentence transformers for domain literature analysis.
    % \item Delivered a production platform, presented to DOE program sponsors.
    % \item Developed the unified UI using Flutter, deployed to production at \url{https://m1.recoil.ise.utk.edu}.
\end{highlights}

\vspace{0.3cm}
\noindent\textbf{SmartShots} - Cross-Platform Application to Improve Childhood Vaccination Rates \hfill Dec 2023 - Dec 2024

\noindent\textbf{Role:} Full-Stack Developer. Built a cross-platform mobile application to improve childhood vaccination tracking and access for families across Tennessee.

\noindent\textbf{Funding:} Tennessee Department of Health.

\begin{highlights}
    \item Developed a scalable Laravel backend and cross-platform Dart/Flutter frontend with real-time data updates, notifications, and alerts, published to the App Store and Google Play.
    \item Integrated community health information for real-time access to nearby vaccination providers, collaborated with the Tennessee Department of Health.
    \item \textit{Winner, IISE DAIS Student Mobile App Competition 2024.}
    \item \url{https://play.google.com/store/apps/details?id=com.ilab.smartshots}
    \item \url{https://apps.apple.com/us/app/smartshots-tn/id6526502640}
\end{highlights}

\vspace{0.3cm}
\noindent\textbf{Active Caregiver's Toolkit (ACTAPP)} - Mobile App for Physical Activity \hfill Jul 2024 - Present

\noindent\textbf{Role:} Mobile Developer. Developed a mobile health application targeting cardiovascular disease risk reduction through physical activity interventions for rural Appalachian caregivers.

\noindent\textbf{Funding:} Hillman Emergent Innovation (HEI).

\begin{highlights}
    \item Implemented Dart 3.0, Material Design, and state management using GetX for a state management.
    \item Integrated a ChatBot for FAQ answering and page\/interaction tracking metrics for data analysis.
    \item \url{https://testflight.apple.com/join/PehxA8aW}
    \item \url{https://play.google.com/store/apps/details?id=com.ilabutk.fe_rich}
\end{highlights}

% \vspace{0.3cm}
% \noindent\textbf{Pathology Cancer Staging - UT Medical Center} \hfill Jan 2025-Present
% \noindent Development of Flutter UI and web map services for interactive demographic cancer visualization.


% ----------------------------
\section{HONORS AND AWARDS}
% ----------------------------
\small

\begin{itemize}[topsep=0pt, parsep=0pt, partopsep=0pt, itemsep=3pt, leftmargin=10pt]
    \item \textbf{IISE Future Faculty Fellows} (2025--2026), USA -- Selected
    \item \textbf{Best Paper Award}, CIE51 (51\textsuperscript{st} Conference on Computers and Industrial Engineering), Sydney, Australia (2024). \textit{``Emerging AI and Cognitive Digital Twin Technologies Towards Low-Carbon Multimodal Freight Transport Systems''}
    \item \textbf{IISE DAIS Student Mobile App Competition -- Winner}, SmartShots Project, Montreal (2024)
    \item \textbf{HIDA Helmholtz Visiting Researcher}, Karlsruhe Institute of Technology (KIT), Germany -- Awarded not taken (2024)
    \item \textbf{Holiday Fellowship}, University of Tennessee, Knoxville (2022, 2023, 2024)
\end{itemize}

% ----------------------------
\section{SELECTED PUBLICATIONS}
% ----------------------------
\small

\noindent TH Wyatt, S Lowe, J Tupayachi, X Li, CA McNeely, X Wang, PD Taylor et al., \textbf{Design and usability testing of SmartSHOTS: A mobile app to reduce vaccine barriers for children 0-24 months}, \textit{Health Informatics Journal 32(1)} (2026).

\noindent J Tupayachi, AN Khan, X Li, \textbf{Scalable decentralized prognostics for industrial systems under data heterogeneity}, \textit{Computers and Electrical Engineering 133, 111023} (2026).

\noindent J Tupayachi, H Xu, OA Omitaomu, MC Camur, A Sharmin, X Li, \textbf{Towards next-generation urban decision support systems through AI-powered construction of scientific ontology using large language models}, \textit{Smart Cities 7(5), 2392-2421} (2024).

\noindent J Tupayachi, MM Ferguson, X Li, \textbf{A Simulation-Based Real-Time Deep Reinforcement Learning Approach for Fighting Wildfires}, \textit{ANNSIM 2024, 1-12} (2024).

\noindent H Xu, X Li, J Tupayachi, JJ Lian, OA Omitaomu, \textbf{Automating Bibliometric Analysis with Sentence Transformers and RAG}, \textit{Proc.\ 2nd ACM SIGSPATIAL Workshop on Advances in Urban-AI} (2024).

\noindent X Li, J Tupayachi, A Sharmin, M Martinez Ferguson, \textbf{Drone-aided delivery methods, challenges, and the future: A methodological review}, \textit{Drones 7(3), 191} (2023).

\noindent J Tupayachi, L Silva, \textbf{Better Efficiency on Non-performing Loans Debt Recovery and Portfolio Valuation Using Machine Learning Techniques}, \textit{POMS Lima, Peru} (2022).


\vspace{0.2cm}
\noindent\textbf{Pre-Prints}

\noindent X Li, H Xu, J Tupayachi, O Omitaomu, X Wang, \textbf{Empowering Cognitive Digital Twins with Generative Foundation Models: Developing a Low-Carbon Integrated Freight Transportation System}, \textit{arXiv:2410.18089} (2024).

\noindent J Tupayachi, MC Camur, K Heaslip, X Li, \textbf{Predicting Electric Vehicle Charging Station Demand: Spatial-Temporal GNNs with Integrated Weather and Traffic Data}, \textit{arXiv:2510.09048} (2024).

\noindent J Tupayachi, X-Y Yu, X Li, Z Liu, K Birdwell, \textbf{Explainable Machine Learning for Extreme Weather Event Prediction at the Oak Ridge Reservation}, \textit{SSRN 6409258} (2025).

\noindent H Xu, J Tupayachi, X-Y Yu, \textbf{Context-Aware Visual Prompting: Automating Geospatial Web Dashboards with Large Language Models and Agent Self-Validation for Decision Support}, \textit{arXiv:2511.20656} (2025).

\noindent J Tupayachi et al., \textbf{Automated End-to-End Geospatial Analysis: An AI-Driven Interface}, \textit{arXiv} (2025).


\vspace{0.2cm}
\noindent\textbf{Submitted Papers}

\noindent H Xu, Y Sun, J Tupayachi, O Omitaomu, S Zlatanova, X Li, \textbf{Towards Autonomous Urban Logistics Optimization via Generative AI and Agentic Digital Twins}, \textit{Computers \& Industrial Engineering} (ID: CAIE-D-25-05889).

\noindent J Tupayachi et al., \textbf{CliniSense AI for Automated Clinical Skills Assessment with Real-Time Feedback}, \textit{npj Digital Medicine}.

\noindent J Tupayachi, MC Camur, K Heaslip, X Li, \textbf{Spatio-Temporal GCNs for EV Charging Demand Forecasting}, \textit{Applied Energy} (ID: 30181).


\vspace{0.2cm}
\noindent\textbf{Upcoming Publications}\\
\noindent Automated Control: The Tragedy of the Commons in End-to-End AI-Driven Systems

\noindent Literature Survey for New Materials Discovery Powered by AI Agents

\noindent Vision Transformer-Based Quality Assurance Engine for Wind Speed and Direction Classification

% (moved to Pre-Prints)




\vspace{0.3em}
\section{TECHNICAL SKILLS}
\noindent\textbf{Languages:} Python, SQL, Scala, Java, R, SAS, STATA, SPSS, Bash, Dart, PHP, C++, CUDA.
\textbf{AI Coding Tools:} GitHub Copilot, Claude Code, Gemini.
\textbf{Cloud:} AWS (S3, EMR, Redshift), Microsoft Azure, Google Cloud.
\textbf{Distributed \& Streaming:} Spark, Hadoop/MapReduce, Hive, Kafka, EMR, batch and real-time streaming ETL/ELT.
\textbf{Data Engineering:} Delta Lake, Databricks Workflows, Unity Catalog, dbt, Airflow, Prefect, data modeling, entity resolution, data quality, profiling, anomaly detection, and monitoring/alerting.
\textbf{Data Warehousing:} Snowflake, Amazon Redshift, dimensional modeling, lakehouse and data-warehouse architecture, query and performance optimization.
\textbf{Databases:} PostgreSQL, MySQL, SQL Server, NoSQL: MongoDB, Cassandra.
\textbf{Analytics \& ML:} Statistical modeling, predictive modeling, clustering, regression, tree-based and boosting models (Random Forest, LightGBM, XGBoost), CNNs, LLMs, RAG, and agentic AI workflows.
\textbf{Frameworks:} Django, Flask, TensorFlow, PyTorch, Flutter, Laravel.
\textbf{DevOps \& Practices:} UNIX/Linux, Docker, Git, Jenkins, CI/CD, Agile/Scrum, KVM, Slurm.
\textbf{Optimization:} Gurobi, CPLEX, AnyLogic, NetworkX, model tuning and calibration.
\textbf{Integration:} REST and GraphQL APIs, data-privacy-aware design.
\textbf{GIS:} WMS, GeoServer, OverPass API.


% ----------------------------
\section{SYNERGISTIC ACTIVITIES}
% ----------------------------
\small

\noindent\textbf{Community Service.}
\textbf{University of Tennessee Graduate Student Senate} - Senator, Industrial and Systems Engineering, 2025. Represented students, advocated for academic and professional development opportunities, and fostered graduate student engagement.

\vspace{0.3em}
\noindent\textbf{Journal \& Conference Reviewer.}
PLOS ONE, Information Systems Frontiers (Springer), IUI 2026 (IEEE/ACM), Journal of Asian Architecture and Building Engineering (Elsevier), IEEE Internet of Things Magazine, Software: Practice and Experience (Wiley), IISE Transactions (Taylor \& Francis), Optimization Letters (Springer Nature), Journal of Robotics (Wiley), Journal of Sensors (Wiley), Winter Simulation Conference 2025, IISE Annual Conference 2025, IEOM 2025.

\vspace{0.3em}
\noindent\textbf{Conference Presenter.}
Federated Learning Fault Detection: Towards a Decentralized ML Framework - IISE 2023.
Rule-based Automated Cancer Staging from Scanned Pathology Reports - IISE 2024.
Empowering Simulation Modeling: An Automated Ontology Framework Enhanced by LLMs - INFORMS 2024.
Conversational Geographic Question Answering: LLMs \& Continuous RAG - SIGSPATIAL 2024.
Automated Literature Review for Solid Lithium Breeder - ANS Annual Meeting 2025, Denver, Colorado.


\end{document}




